\subsubsection*{Solution}
\paragraph*{Step-by-Step Derivation}

\subparagraph*{1. Represent the driving field}

Let $\sigma=+1$ denote left-circular polarization and $\sigma=-1$ right-circular polarization. The three pulses therefore have

\[
(\ell_1,\sigma_1)=(-1,+1),\qquad (\ell_2,\sigma_2)=(2,-1),\qquad (\ell_3,\sigma_3)=(1,+1).
\]

Their positive-frequency field is schematically

\[
\mathbf E^{(+)}(\rho,\phi,t) = \sum_{j=1}^{3} E_j(\rho)f(t-t_j) e^{i\ell_j\phi} \mathbf e_{\sigma_j} e^{-i\omega_0t+i\varphi_j},
\]

where $t_j=(0,30,60)\,\mathrm{fs}$ and $\varphi_j$ are the relative carrier phases.

Because this sum contains different $\ell_j$ and $\sigma_j$, it is not an eigenstate of either OAM or helicity during temporal overlap.

\subparagraph*{2. The pulses overlap significantly}

Assuming the quoted $50\,\mathrm{fs}$ duration is the intensity FWHM of a Gaussian pulse,

\[
\frac{I_j(t)}{I_0} = \exp\left[ -4\ln 2\left(\frac{t-t_j}{50\,\mathrm{fs}}\right)^2 \right].
\]

At a separation of $30\,\mathrm{fs}$,

\[
\frac{I_j(t_j\pm30\,\mathrm{fs})}{I_0} = \exp\left[-4\ln2\left(\frac{30}{50}\right)^2\right] = 0.368567304323.
\]

Thus adjacent pulses retain about $36.9\%$ of their peak intensity at one another's peak. They cannot generally be treated as independent. Time-delayed vortex drivers are known to produce harmonics with time-dependent or distributed OAM rather than one fixed topological charge; this is the basis of HHG self-torque (\href{\detokenize{https://pubs.acs.org/doi/10.1021/acsphotonics.4c01320}}{ACS Photonics}).

\subparagraph*{3. Allowed photon channels for harmonic 23}

Let $n_j$ be the number of fundamental photons drawn from pulse $j$ in a particular lowest-order photon channel. Since all pulses have the same frequency $\omega_0$,

\[
n_1+n_2+n_3=23.
\]

For an isotropic gas returning to the same internal state, spin and orbital angular momenta of the electromagnetic field are conserved separately (\href{\detokenize{https://pmc.ncbi.nlm.nih.gov/articles/PMC11317519/}}{Nature Communications}). Hence

\[
\ell_{23}=n_1\ell_1+n_2\ell_2+n_3\ell_3 =-n_1+2n_2+n_3,
\]

and

\[
\sigma_{23} =n_1\sigma_1+n_2\sigma_2+n_3\sigma_3 =n_1-n_2+n_3.
\]

Because the emitted photon must have helicity $\sigma_{23}=\pm1$,

\[
n_1-n_2+n_3=\pm1.
\]

\subparagraph*{Left-helicity harmonic: $\sigma_{23}=+1$}

Solving

\[
n_1+n_2+n_3=23,\qquad n_1-n_2+n_3=1
\]

gives

\[
n_2=11,\qquad n_1+n_3=12.
\]

Writing $n_1=k$ and $n_3=12-k$,

\[
\ell_{23} =-k+2(11)+(12-k) =34-2k, \qquad k=0,1,\ldots,12.
\]

Therefore,

\[
{ \sigma_{23}=+1,\qquad \ell_{23}\in\{10,12,14,\ldots,34\}. }
\]

\subparagraph*{Right-helicity harmonic: $\sigma_{23}=-1$}

Similarly,

\[
n_2=12,\qquad n_1+n_3=11.
\]

With $n_1=k$ and $n_3=11-k$,

\[
\ell_{23} =-k+2(12)+(11-k) =35-2k, \qquad k=0,1,\ldots,11.
\]

Thus,

\[
{ \sigma_{23}=-1,\qquad \ell_{23}\in\{13,15,17,\ldots,35\}. }
\]

The strength and phase of each mode depend on the pulse intensities, relative carrier phases, radial profiles, gas species, microscopic dipole response, and macroscopic phase matching. Mixing vortex modes in gas HHG generally produces an OAM spectrum whose weights depend on precisely these quantities (\href{\detokenize{https://journals.aps.org/pra/abstract/10.1103/PhysRevA.105.023107}}{Physical Review A}).

Moreover, isolated single-color circular pulses do not efficiently drive ordinary recollision-based gas HHG, because circular polarization prevents electron return (\href{\detokenize{https://journals.aps.org/pra/abstract/10.1103/PhysRevA.48.R3437}}{Physical Review A}, \href{\detokenize{https://link.aps.org/accepted/10.1103/PhysRevLett.119.063201}}{Physical Review Letters}). The relevant emission here would arise primarily where the counter-rotating components overlap and produce a locally noncircular field.
